
\documentclass[12pt,a4paper]{article}
\usepackage[utf8]{inputenc}
\usepackage[T1]{fontenc}
\usepackage{graphicx}
\usepackage{geometry}
\usepackage{array}
\usepackage{xcolor}

\begin{document}
\begin{titlepage}
    \begin{table}[!h]
		\begin{center}
			\resizebox{\linewidth}{!} {
				\begin{tabular}{c c}
					\bsc{\LARGE{\textbf{UNIVERSITE DE YAOUNDÉ I}}}	& \bsc{\LARGE{\textbf{UIVERSITY OF YAOUNDE 1}}}\\ \\ ** & **\\\\ 
					\bsc{\LARGE{\textbf{ECOLE NATIONALE SUPÉRIEURE }}}	& \bsc{\LARGE{\textbf{NATIONAL ADVANCED SCHOOL OF}}}\\ \\ 
					\bsc{\LARGE{\textbf{POLYTECHNIQUE DE YAOUNDE}}}			& \bsc{\LARGE{\textbf{ENGINEERING OF YAOUNDE}}}\\ \\ ** & **\\  \\
					\bsc{\LARGE{\textbf{DEPARTEMENT DE GÉNIE}}}		& \bsc{\LARGE{\textbf{DEPARTEMENT OF COMPUTER}}}\\ \\
					\bsc{\LARGE{\textbf{INFORMATIQUE}}}			& \bsc{\LARGE{\textbf{ENGINEERING}}}\\ \\** & **
				\end{tabular}
			}
		\end{center}
	\end{table}

	\begin{figure}[!h]
		\begin{center}
			\includegraphics[scale=0.5]{po.png} 
		\end{center}
	\end{figure}
    

  
 	\begin{center}
				\large{\textbf{\bsc{ INTRODUCTION AUX TECHNIQUES D'INVESTIGATION NUMERIQUE}}}}}}}
			\end{center}
	}} 
	%\vspace{0.2cm}{INTRODUCTION AUX TECHNIQUES D'INVESTIGATION NUMERIQUE}}}}
  
  \vspace{15px}
	\colorbox{gray}{\parbox{\linewidth}{
			\begin{center}
				\Large {\textbf{\bsc{Résumé du cours}}}
			\end{center}
	}} 
	\vspace{1cm}
    
    	Présenté par:\\
	\vspace{8px}
	\textbf{NGWAMBE BEKWADI MARIELLA M. - 22P040}\\
	\vspace{10px}
    	  Sous la supervision de:\\
	\vspace{8px}
	\textbf{M. MINKA THIERRY}
	\vspace{30px}
    
	\centering {Année académique \oldstylenums{2025} - \oldstylenums{2026}}\\ 
	\vspace{6px}
\end{titlepage}
 \chapter \huge{{\textbf{{RESUME}}
 \normalsize
 \paragraph{}
 L’investigation numérique n’est pas seulement une pratique technique : elle cons-titue aujourd’hui une discipline philosophique et scientifique, qui interroge notre rapport à la vérité, à la preuve et à la justice dans un monde digitalisé. Dès lors, l’investigateur numérique ne se limite pas à un technicien, il devient philosophe-praticien, chargé d’exhumer, préserver et interpréter les traces numériques qui façonnent la mémoire collective. À l’ère de la société numérique, l’existence humaine se déploie simultanément dans deux dimensions : physique et digitale. Ce « double numérique » engendre de nouvelles réalités :
\begin{itemize}
    \item une ontologie numérique, où l’être existe à travers ses traces  di-gitales
    \item une phénoménologie des données, où chaque action laisse une empreinte,
    \item et une métaphysique digitale, qui redéfinit nos modes de relation et de confiance.
\end{itemize}\item 
Mais cette digitalité s’accompagne de tensions.La quête de transparence se heurte très souvent au droit fondamental à l’intimité,plaçant l’investigateur au carrefour d’un dilemme : protéger la vérité tout en respectant la vie privée.La preuve numérique transforme aussi notre rapport à la vérité. Contrairement à la preuve matérielle, stable et tangible, la preuve digitale est volatile, mutable et fragile. Son authenticité ne repose plus sur l’objet mais sur une chaîne de confiance et des mécanismes cryptographiques. À cette mutation s’ajoute une crise de la vérité : manipulations algorithmiques (deepfakes), multiplication des sources et fragmentation du réel compliquant la tâche de l’investigateur, désormais archiviste du monde numérique. 
\paragraph{}
Pour analyser ces réalités, l’investigation numérique s’appuie sur des fondements mathématiques :
\begin{itemize}
    \item La théorie de l’information (Shannon) pour mesurer l’incertitude et détecter des anomalies.
    \item La théorie des graphes, qui modélise les interactions sociales et techniques.
    \item La théorie du chaos, rappelant la sensibilité extrême des systèmes numériques aux moindres variations.
\end{itemize}
À cela s’ajoute un élémént majeur:la révolution quantique, qui bouleverse les paradigmes traditionnels. L’investigateur doit désormais composer avec la superposition, l’intrication et l’incertitude, remettant en cause les certitudes binaires de la preuve classique. Cette révolution s'accompagne des algorithmes de Shor et Grover menaçant les systèmes cryptographiques classiques. 
Par la suite se présente le paradoxe de l’authenticité invisible illustrant les tensions de notre époque : plus une preuve est authentique et vérifiable, plus elle risque de compromettre la confidentialité ; inversement, plus elle reste protégée, plus son authenticité devient difficile à établir. Les protocoles cryptographiques de type Zero-Knowledge Non-Repudiation (ZK-NR) offrent une solution en permettant de vérifier sans révéler, redéfinissant ainsi la preuve comme processus plutôt qu’objet.

\paragraph{}Dans le sens des obligations vient l'éthique et la responsabilité de l’investigateur. L’investigateur numérique doit constamment arbitrer entre trois tensions éthiques: la transparence vs vie privée, l'efficacité vs proportionnalité,l'innovation vs responsabilité.Sa mission s’articule autour d’une charte claire : préserver l’intégrité des preuves, respecter la dignité numérique des individus, reconnaître les limites de sa connaissance et toujours rechercher la vérité au-delà de la simple conviction.La trace numérique, loin d’être une donnée brute, est une manifestation d’existence : elle révèle des intentions, s’inscrit dans une temporalité plurielle et nécessite une interprétation fine des données. Or, dans un monde où la mémoire est quasi-parfaite, l’éthique post-quantique impose de préserver la possibilité de l’oubli et de garantir la justice même lorsque la vérité est « cryptée ». L'un des apports les plus remarquables de l'investigation numérique est la mise en place d'une charte déontologique pour l'investigateur. Elle s'incarne par les dix commandements(en analogie aux 10 commandements chrétiens) qui rappellent que la technique doit rester au service de la justice et du respect des droits humains. Les piliers fondamentaux de la pratique éthique sont l'intégrité, la proportionnalité la responsabilité et le service. L'investigateur doit être minutieux,rigoureux et capable de résister aux pressions contraires à l'éthique.

 \paragraph{}
Par ailleurs,il est important de faire un bref historique de l’investigation numérique. L’évolution de la discipline révèle une professionnalisation progressive :
\begin{itemize}
    \item 1990-2000 : L’Opération Sundevil et l’arrestation de Kevin Mitnick marquent la naissance de standards comme la chaîne de conservation.
    \item 2000-2010 : Affaires Enron et McKinnon conduisent à la standardisation internationale (ISO/IEC 27037).
    \item 2010-2020 : Big Data, cloud et blockchain imposent de nouvelles méthodes (Silk Road, Panama Papers).
    \item 2020-aujourd’hui : L’ère post-quantique et l’IA (attaque SolarWinds) redéfinissent les techniques d’attribution et d’analyse.

\end{itemize}
Les grandes affaires de Stuxnet à WannaCry illustrent l’importance des métadonnées, de l’ingénierie inverse, de l’analyse comportementale et des méthodes automatisées dans la lutte contre des menaces toujours plus sophistiquées.L'investigation numérique est circonscrite par un cadre théorique et des modèles pratiques.Inspirée du principe de Locard (« toute action laisse une trace »), l’investigation numérique distingue les traces primaires (logs, fichiers temporaires, registres) et les traces secondaires (métadonnées, flux réseau, empreintes comportementales). Pour structurer la démarche, plusieurs modèles théoriques existent : Le DFRWS (identification, préservation, collection, analyse, présentation),le modèle de Casey, qui intègre préparation, scène physique et scène numérique,la norme ISO/IEC 27037 qui formalise l'identification,l'acquisition et la documentation des preuves.Ces modèles sont renforcés par l’usage appliqué des mathématiques : l'entropie pour détecter anomalies et stéganographie,la distance de Hamming pour identifier des variantes de malware, ou encore graphes pour analyser réseaux sociaux et flux de données. L'investigation numérique est porteuse de nombreux paradigmes et de perspectives en étroite évolution avec les technologies actuelles et futures dont quelques uns ont retenu notre attention:
\begin{itemize}
    \item Le Digital Forensics as a Service (DFaaS), automatisé et scalable grâce au cloud.
    \item Une approche proactive, intégrant le threat hunting et la collecte anticipée.
    \item Demain, l’investigation sera marquée par l’adaptation post-quantique, l’usage massif de l’intelligence artificielle et l’intégration de la cryptographie avancée dans les processus de vérification.
\end{itemize}
\paragraph{}
Les normes et standards internationaux constituent aujourd’hui un pilier incontournable de la criminalistique numérique,car ils garantissent que les preuves collectées soient fiables, exploitables et recevables devant les juridictions. L’ISO/IEC 27037:2012 en pose les bases en encadrant l’identification, la collecte, l’acquisition et la préservation des éléments numériques, avec des principes fondamentaux comme la pertinence, la fiabilité, la suffisance et la documentation complète. Une fois la preuve préservée, la question de la validité des méthodes se pose:c’est l’objet de l’ISO/IEC 27041:2015, qui veille à ce que les techniques employées, qu’il s’agisse de live forensics (RFC 3227), de dead forensics (NIST SP 800-86), d’investigations réseau (IETF) ou mobile (NIST SP 800-101), soient reconnues et reproductibles. Mais la valeur d’une preuve ne réside pas uniquement dans sa collecte : elle doit aussi être analysée et interprétée de manière rigoureuse. C’est précisément ce que formalise l’ISO/IEC 27042:2015 à travers un cycle structuré allant de la préparation à l’extraction, puis de l’analyse à l’interprétation et enfin à la rédaction du rapport.
 Enfin, pour offrir une vision globale et cohérente du processus d’investigation, l’ISO/IEC 27043:2015 propose un modèle intégrant toutes les étapes : de la préparation et la détection jusqu’à la présentation des résultats et le retour d’expérience. Ensemble, ces normes assurent une méthode structurée,harmonisée et crédible de l’investigation numérique, favorisant non seulement la solidité des enquêtes mais aussi la coopération internationale face à une cybercriminalité en constante évolution.L'intégration cruciale de la forensique explicable (XAI) pour rendre les analyses IA juridiquement admissibles, la nécessité de frameworks de résilience anti-forensique proactifs et adaptatifs (optimisés par le Trilemme CRO et les preuves ZK-NR) face à l'intensification des techniques d'obfuscation, et réalise un benchmarking mondial des pratiques forensiques. Ce dernier compare l'excellence technique du FBI/NIST (normalisation, évaluation CRO), la rigueur procédurale de Scotland Yard (principes ACPO, cryptographie post-quantique, blockchain/ZK-NR pour l'audit), et la précision métrologique du BKA allemand (standards techniques, validation d'outils). L'ensemble souligne l'évolution vers une "forensique inviolable" où l'investigateur doit être un "gardien de l'intégrité numérique", préparé à l'ère post-quantique et capable d'intégrer intelligence artificielle, cryptographie avancée et coopération internationale pour protéger la vérité face à toutes les manipulations.

L'investigation numérique fait immerger des pratiques avancées et des techniques modernes. Ici,on s'appesantit sur les pratiques opérationnelles comme la mise en place d'un laboratoire forensique, les procédures standardisées, les formations continues et les veilles technologiques. Les techniques de pointe incluant la memory forensics, la timeline analysis et les environnements virtuels ainsi que l'analyse moderne(6G/5G, protocoles chiffrés). Un autre pan important est la confrontation entre forensique et anti-forensique, les techniques de dissimulation, de stéganographie et l'implication de l'intelligence artificielle qui joue un rôle double-face notamment attaquer et défendre.Ces principes,ces techniques et ces visions sont encadrés par des éléments juridiques appliqués dans les cadres régionaux:
\begin{itemize}
    \item Etats-Unis:Féderal Rules  of Evidence, CFAA, Stored Communications Act.
    \item Europe :Règlement eIDAS, RGPD, Convention de Budapest.
    \item Convention de Malabo et au Cameroun les lois 2010/012 et 2010/013 constituent la base légale de l’investigation numérique, renforcée récemment par la loi 2024/017.Il est nécessaire de former des experts agréés, de développer une jurisprudence locale et de résoudre les défis de compétence technique et juridique.
\end{itemize}
Des cas emblématiques tels que BTK Killer,Stuxnet et un cas fictif CyberFinance Cameroun 2025 montre clairement la gestion d'une attaque par ransomware,depuis la détection jusqu'à la remédiation en passant par la collecte des preuves, l'analyse forensique,l'attribution de l'attaque et la préparation judiciaire.    
\paragraph{}
L’investigation numérique se situe au croisement de la technique, de la science, de l’éthique et de la philosophie. Elle dépasse le simple cadre judiciaire pour devenir une véritable praxis de liberté, garantissant la mémoire collective et la justice dans un monde où le réel est fragmenté et manipulable. Ainsi, l’investigateur moderne est bien plus qu’un technicien : il est gardien de l’intégrité informationnelle, médiateur entre transparence et intimité, et acteur essentiel de la confiance numérique.
\end{document}